% ========================================
% Document Class and Package Setup
% ========================================
\documentclass[reprint,amsmath,amssymb,aps,]{revtex4-1}

\usepackage{graphicx}% Include figure files
\usepackage{dcolumn}% Align table columns on decimal point
\usepackage{bm}% bold math
\usepackage[colorlinks=true,linkcolor=blue,filecolor=blue, urlcolor=blue]{hyperref}
\usepackage{braket}
\usepackage{caption}
\usepackage{threeparttable}
\numberwithin{equation}{section}% (p.37)


% ========================================
% Document Begin
% ========================================
\begin{document}

\captionsetup[figure]{justification=raggedright,
    labelfont={bf},name={Fig.},labelsep=period}
\captionsetup[table]{
    labelfont={bf},name={Table.},labelsep=period}


% ========================================
% Title, Author, and Date
% ========================================
    \title{Transition Theory of An Electron \\ Traveling from Uncertain to Causal Basis}
    \author{Satoshi Hanamura}
    \email{hana.tensor@gmail.com}
    \date{September 21, 2019}
% ========================================
% Abstract
% ========================================
        \begin{abstract}
        One of Einstein's most famous quotes is; `I am at all events convinced that God does not play dice'. This study attempted to be a conversion of ``uncertainty principle'' to a ``certainty principle.'' With the previous electron model, bare electrons moved discretely, and the photons surrounding them moved continuously. In this study, we shall notice the traveling of free electrons with thermal conductance, and how to determine the unique discrete traveling point by using thermal potential energy gradient.
        \end{abstract}
    \maketitle

% = = * * = =  * * = = * * = = * * = = * * = = * * = = * * = = * *
% ========================================
% Section 1: INTRODUCTION
% ========================================
\section{\label{sec:INTRODUCTION}INTRODUCTION}


    The thermal conductance of single-molecule junctions has been fully characterized \cite{Pauly} recently. This study was also confirmed that the fundamental law of heat conduction in classical physics, Wiedemann–Franz law was established at the atomic level.

    This law is valid in metals. This is because free electrons carry most of both thermal conduction and electric conduction.
    
    In this study, we reflected how free electrons will travel with its oscillation using the previous electron model \cite{Hanamura1}. The feature of the previous electron model is that bare electrons move discretely. In this study, we will proceed with discussions on the assumption of discretely moving free electrons.

    The movement of free electrons should be treated as thermal conduction, not as radiation. This is because thermal potential energy possessed by bare electrons is in an oscillatory relationship by radiation and absorption. This means that the discrete movement of bare electrons is slower than the speed of light. Therefore, the photons that wrap the bare electrons move slower than the speed of light.

    
    The previous electron model has refined the concept of virtual photons. There are two adjustments. One is that virtual photons that have been thought to occur only during interaction are assumed to exist continuously. The other adjustment is that a virtual photon we have conceived is not a virtual but a real entity. It could be said that the photon energy is confined by the central force of the bare electrons, which would be a kind of black hole. Therefore, the idea that an electron has internal structure which is consisted of a real photon and bare electrons has not come out. It is emphasized again that this can be said to be a state in which photons have mass.


        \begin{figure}[h!]
            \centering 
            \includegraphics[scale=1]{f1_random.png}
                \caption{A schematic of the manner in which a virtual photon can be moved as a simple harmonic oscillator with the emergence and disappearance of bare electrons at the fixed points $x=a$ and $x=-a$. $(\bm{\mathrm{a}}) \ \theta = 0 $,  $(\bm{\mathrm{b}}) \ \theta = \frac{\pi}{2} $,  $(\bm{\mathrm{c}}) \ \theta = \pi $,  $(\bm{\mathrm{d}}) \ \theta = \frac{3\pi}{2} $,  $(\bm{\mathrm{e}}) \ \theta = 2\pi $, where $\theta$ is  electron's phase.} 
            \label{f1_random.png}
        \end{figure}
            % The two bare electrons are two spinors 1 and 2 represented by the blue and green dots, respectively. The open circles express the thermal potential energy, which has a zero value. The arrows within the yellow circles (b) and (d) have two meanings: the direction in which the virtual photon is moving, and the direction to which the thermal potential energy is radiated between the two spinor particles (i.e., the blue and green dots)

        \begin{figure}[h!]
            \centering 
            \includegraphics[scale=1]{f2_random.png}
                \caption{$(\bm{\mathrm{a}})$ $p = \frac{1}{2}$. The bare electron at the origin moves to the left. $(\bm{\mathrm{b}})$ $q = \frac{1}{2}$. The bare electron at the origin moves to the right. The yellow dashed line is a real photon as re-imaged virtual photon. (Steps to the right are performed with probability $p$ and steps to the left with probability $q = 1 - p$.)} 
            \label{f2_random.png}
        \end{figure}


% = = * * = =  * * = = * * = = * * = = * * = = * * = = * * = = * * 
% ========================================
% Section 2: METHODS
% ========================================
\section{\label{sec:METHODS}METHODS}

% ----------------------------------------
% Subsection 2.1:
% ----------------------------------------
\subsection{
    \label{sec:Randwom walk}Random walk}

    In the previous new electron model, we explained that bare electrons move discretely \cite{Hanamura2}.
    
    Since divergence cannot be differentiated at zero-dimensional point, M-theory has attempted to solve the problem of divergence by creating a spatial spread and devising string-like particles.

    The electronic model I devised placed two points instead of strings, and was trying to solve the problem of divergence by connecting these two points to form a one-dimensional line and using the length of the space constructed by the two dots.

    Blue and green dots are equivalent to “bare electrons” (Fig. \ref{f1_random.png}). Each is a black body, and initially the energy of the blue dot is transferred to the green dot. It was assumed that the bare electron attracts the photon by centering force, the force disappears if the thermal energy is radiated.

    The position of the central force that attracts the photon \cite{Hanamura1}, similar phenomena to the center of gravity, the virtual photon appears to swing left and right as a simple oscillator. The bare electron, the blue dot, is not necessary to return to the original location, $o$. The position of the next blue dot can be in a different direction.


        \begin{figure}[h!]
            \centering 
            \includegraphics[scale=1]{f3_random.png}
                \caption{$(\bm{\mathrm{a}})$ If another thermal particle occupies the position of $+a$, the blue dot bare electrons cannot move to $+a $. $(\bm{\mathrm{b}})$ $ p = 1 $ ; By applying the second law of thermodynamics, the blue dot moves to $-a$. Blue bare electrons are generated at coordinate $-a$, which is a green dot. In this situation, the probability of $p$ is 1.} 
            \label{f3_random.png}
        \end{figure}


    Suppose that there is one free electron on the $x$-axis indicating the position of the space. These electrons can travel in both positive and negative directions within the range of virtual photons that enclose bare electrons. The probability of going to is the same.

    By applying this assumption, an electron can move randomly (Figs. \ref{f2_random.png} and \ref{f3_random.png}).


% = = * * = =  * * = = * * = = * * = = * * = = * * = = * * = = * * 
% ========================================
% Section 3: DISCUSSION
% ========================================
\section{\label{sec:DISCUSSION}DISCUSSION}

% ----------------------------------------
% Subsection 3.1:
% ----------------------------------------
\subsection{
    \label{sec:Applying the second law of thermodynamics}Applying the second law of thermodynamics}


    The discussion here implies the key to unraveling why the Seebeck effect occurs from a quantum mechanical point of view.

    As seen in the previous chapter, a complete random walk occurs when there are no temperature gradients or potential gradients on the left and right. In this section, we examine how electrons move in the presence of a temperature gradient or potential gradient.

    The state where electrons are placed at the origin of the $x$-axis indicating the position of the one-dimensional space is the same as in the previous section. The difference is that there is another electron at coordinate A on the $x$-axis. In this case, the electron at the origin tends to advance in the lower temperature gradient, that is, in the negative $x$-axis direction. This is based on the second law of thermodynamics.


        \begin{figure}[h!]
            \centering 
            \includegraphics[scale=1]{f4_random.png}
                \caption{$(\bm{\mathrm{a}})$ A grid extended from one dimension to two dimensions. $(\bm{\mathrm{b}})$ Random walker basis; Another electron exists at point A. The potential energy of bare electrons moves from the origin to point B, C, or D. $(\bm{\mathrm{c}})$ Causal basis; A grid assumes that three electrons are occupied at fixed positions A, B, and D. The bare electron located at the origin (blue dot) is forced to travel to point C having the lowest thermal potential energy gradient. }
            \label{f4_random.png}
        \end{figure}


% ----------------------------------------
% Subsection 3.2:
% ----------------------------------------
\subsection{
    \label{Extension to two dimensions}Extension to two dimensions}

    Fig. \ref{f4_random.png} shows a lattice in which the electron movement range is expanded from one dimension to two dimensions. The distances from the origin to point A, B, C, D are equal, $r$. In Fig. \ref{f4_random.png} (b), the probabilities that the electrons at the origin move to points B, C, and D are equal, and all $ p = \frac{1}{3} $.

    This is because point A has the highest energy in energy gradient among points A, B, C, and D. The electron located at point O cannot continue to be at the same point O. Bare electron at point O oscillates and forces to be travel to another position.

    In this state, it can be said that electrons move randomly.


% ----------------------------------------
% Subsection 3.3:
% ----------------------------------------
\subsection{
    \label{sec:No random events}No random events}

    Now assume that there are three electrons at points A, B and D (Fig. \ref{f4_random.png} (c)). At this time, the energy gradient of the thermal potential energy is equal at points A, B, and D. On the other hand, point C has the lowest value.

    Therefore, thermal potential energy of the bare electron  located at the origin must travel to point C. The probability of moving to point C is $p = 1$, which is no longer a random walk.
    
    In the previous electron model \cite{Hanamura1}, it is assumed that the bare electrons included in a photon have thermal potential energy, so we have explained this model using thermodynamics.

    This study follows the procedure of using thermodynamics as well as the previous study.
    
    Utilizing the gradient of thermal potential energy provides a consistent explanation for the second law of thermodynamics.


% ----------------------------------------
% Subsection 3.4:
% ----------------------------------------
\subsection{
    \label{sec:}Applicability of the principle of minimum action for photons contained in the electron model}
    

    From this perspective, a motivation of whether the principle of minimum action could be applied to geometric optics arises. However, L.D. Landau and E.M. Lifshitz stated that it is impossible to introduce a function equivalent to a Lagrangian for a ray.
    
    
    \begin{quote}
        Pursuing the analogy, we can establish for geometrical optics a principle analogous to the principle of least action in mechanics. However, it cannot be written in Hamiltonian form as $\delta \int L dt = 0$, since it turns out to be impossible to introduce, for rays, a function analogous to the Lagrangian of a particle. Since the Lagrangian of a particle is related to the Hamiltonian $\mathcal{H}$ by the equation $L = \bm{p} \cdot \delta \mathcal{H} / \delta \bm{p} - \mathcal{H}$, replacing the Hamiltonian $\mathcal{H}$ by the frequency $\omega$ and the momentum by the wave vector $\bm{k}$, we should have to write for the Lagrangian in optics $\bm{k} \cdot \delta \omega / \delta \bm{k}- \omega$. But this expression is equal to zero, since $\omega = c k$. The impossibility of introducing a Lagrangian for rays is also clear directly from the consideration mentioned earlier that the propagation of rays is analogous to the motion of particles with zero mass. \cite{Landau}
    \end{quote}

    Landau explained that it is impossible to introduce a Lagrangian into a ray because the propagation of the ray resembles the motion of a mass-less particle. However, virtual photons as real photons which are not behaving as the ray, covering bare electrons behave as if they have mass in the new electron model. This is because the photon makes an accelerating movement as a simple oscillator as shown in Fig. \ref{f1_random.png}.
    

    Since the new electron model consisting of a real photon with bare electrons is different from the premise of the above quote, there would be room for applying a principle similar to the principle of minimum action to photons that travel at an accelerated rate. Finding a specific Lagrangian would need further study.
    
    Applicability of the principle of minimum action for the photons contained in the electron model and derivation of the Lagrangian would need further study.
    

% = = * * = =  * * = = * * = = * * = = * * = = * * = = * * = = * * 
% ========================================
% Section 4: CONCLUSION
% ========================================
\section{\label{sec:CONCLUSION}CONCLUSION}


    The essence of this study is that free electrons move in the direction of lower thermal potential energy gradient, and the direction of movement is not determined stochastically. It was a remarkable achievement to associate the second law of thermodynamics and quantum mechanics through the thermal potential energy gradient.

    It was interesting to point out the question about the randomness of quantum mechanics. Quantum mechanics is fundamentally about randomness and indeterminacy. This study attempted to be a conversion of ``uncertainty principle'' to a ``certainty principle.''  Einstein replied to a letter from Max Born; `I am at all events convinced that God does not play dice' in December 1926. 

    However, this study uses a lattice model, and it can be applied to an Ising model that introduces the spin concept.
    
    The new electron model reinterpreted a virtual photon as an actual photon \cite{Hanamura1}. The bare electron is regarded as an actual thermal entity and has a mass. The new electron model in which bare electrons included in a photon could be a powerful basis for linking the laws of mechanics and geometric optics.
    
    One of the achievements shown in this study is that the traveling of electrons is also the traveling trajectory of real photons. This is because the electron model that remodels a virtual photon into a real photon is applied to this study, i.e., the traveling of an electron can be the traveling trajectory of a real photon surrounding a bare electron. This creates room for applying Fermat's theorem to the traveling of electrons.
    
    %The principle of minimum action was found by Maupertuis and generalized by Euler and Hamilton. As a result, the principle of minimum action and Hamilton was established. The application of the variational principle to quantum mechanics triggered by this study would be worthy of attention.

    What we will expect following this research is a formal unification of the subject and Fermat's principle because Fermat's principle and the principle of minimum action have been completed around a variational principle.

% = = * * = =  * * = = * * = = * * = = * * = = * * = = * * = = * * 
%Insert space for margin layout.


% = = * * = =  * * = = * * = = * * = = * * = = * * = = * * = = * * 
% ========================================
% Bibliography
% ========================================
\begin{thebibliography}{99}
    \bibitem{Pauly} Cui, L.; Hur, S.; Akbar, Z.A.; Klöckner, J.C.; Jeong, W.; Pauly, F.; Jang, S.-Y.; Reddy, P.; Meyhofer, E. Thermal conductance of single-molecule junctions, \textit{Nature} \textbf{572},  628-633 (2019).
    
    \bibitem{Hanamura1} S. Hanamura, A Model of an Electron Including Two Perfect Black Bodies, Zenodo (2018), \href{https://doi.org/10.5281/zenodo.16759284}{doi:10.5281/zenodo.16759284}. [Originally: \href{http://vixra.org/abs/1811.0312}{viXra:1811.0312}]
    
    \bibitem{Hanamura2} S. Hanamura, A New Representation of Spin Angular Momentum, Zenodo (2018), \href{https://doi.org/10.5281/zenodo.16783727}{doi:10.5281/zenodo.16783727}. [Originally: \href{http://vixra.org/abs/1812.0472}{viXra:1812.0472}]
    
    \bibitem{Landau} L.D. Landau \& E.M. Lifshitz, \href{https://archive.org/details/TheClassicalTheoryOfFields/page/n141}{The Classical Theory of Fields}, p.131, 1971.
    
    \bibitem{Klafter} J. Klafter and I. M. Sokolov, \textit{First Steps in Random Walks}, Oxford University Press, 2011.
    
    \bibitem{Cengel} Yunus A. \c{C}engel and Michael A. Boles, Thermodynamics 8th Ed.,
\end{thebibliography}

% = = * * = =  * * = = * * = = * * = = * * = = * * = = * * = = * * 
% ========================================
% Section 5: Appendix and Quotation
% ========================================
\section{Appendix and Quotation}
\label{Appendix}

% ----------------------------------------
% Subsection 5.1: Random walk and Fourier transform
% ----------------------------------------
\subsection{Random walk and Fourier transform}

    The state of the random walk can be formulated using Fourier transform.
    
    \begin{quote}
        "We consider a particle (a ``random walker'') starting at site 0. Its motion can be described as jumps (steps) to one of the neighboring sites. Different steps of the random walk are considered independent. Steps to the right are performed with probability $p$ and steps to the left with probability $q = 1 - p$. The position of the walker after $n$ steps is determined by the number of steps to the right and to the left. We now introduce an instrument that allows for finding the positions of the walker.
        
        Let us consider the expression $p e^{i \theta} + q e^{i \theta}$. The coefficient preceding $e^{i \theta}$ is the probability that the first step was to the right, and the coefficient preceding $e^{-i \theta}$ is the probability that the first step was to the left. If we now consider the square of this expression, $(p e^{i \theta} + q e^{i \theta} )^2$, and expand it using the binomial formula $(p e^{i \theta} + q e^{i \theta} )^2 = p^{2} e^{2 i \theta} + 2 p q + q^{2} e^{-2 i \theta}$, we see that the coefficient of the first term of the expansion gives the probability that the first two steps were to the right, i.e., that the walker arrives at site $j = -2$. Generalizing these findings to any $n$, we see that the coefficient preceding  $e^{i \theta j}$ in the expansion of the polynomial $( p e^{i \theta} + q e^{-i \theta})^{n}$ gives us the probability $P_{n} (j)$ from the corresponding expression via Fourier transform: \cite{Klafter}"
    \end{quote}


    \begin{equation}
        P_{n} (j) = \frac{1}{2 \pi} \int_{-\pi}^{\pi} \left( p e^{i \theta} + q e^{-i \theta} \right)^{n} e^{-i \theta j} d\theta
        .
    \end{equation}


    If the walker has the same probability of moving left and right, i.e., $p = q = \frac{1}{2}$, the characteristic function is as follows:
    
    \begin{equation}
        p e^{i \theta} + q e^{-i \theta} = \frac{ \left( e^{i \theta} + e^{-i \theta} \right) }{2} = \cos \theta
        .
    \end{equation}


% ----------------------------------------
% Subsection 5.2: Seebeck effect
% ----------------------------------------
\subsection{Seebeck effect}


    The Seebeck effect is one of thermoelectric effects.

    \

    "Consider two wires made from different metals joined at both ends (junctions), forming a closed circuit. Ordinarily, nothing will happen. However, when one of the ends is heated, something interesting happens: A current flows continuously in the circuit. This is  called the Seebeck effect, in honor of Thomas Seebeck, who made this discovery in 1821. The circuit that incorporates both thermal and electrical effects is called a thermoelectric circuit, and a device that operates on this circuit is called a thermoelectric device. \cite{Cengel}"


% ========================================
% End of Document
% ========================================
\end{document}
